% !PRINTED 187
\weberpage{187}

Nun weist man durch Differentiation der Periodengleichungen
(8) nach, daß die Funktion
\[
S(u) = \frac{\partial^2 \sigma}{\partial u^2} + 4\eta_1\frac{\partial \sigma}{\partial \omega_1} + 4\eta_2\frac{\partial \sigma}{\partial \omega_2} + 4ru^2\sigma
\]
den Periodengleichungen (8) genügt, und also gleich $C\sigma(u)$ ist,
worin $C$ von $u$ unabhängig ist. Die Entwickelung von $\sigma(u)$ nach % [sic] Printed sigma(u); the argument requires S(u).
Potenzen von $u$ fängt aber mit $u^3$ an, während $\sigma(u)$ mit $u$ anfängt,
und folglich muß $C=0$ sein.

Nach (7) beginnt die Entwickelung von $\partial\sigma/\partial\omega_1$ und $\partial\sigma/\partial\omega_2$
erst mit $u^5$, und die ersten Glieder von $\partial^2\sigma/\partial u^2$ und $4ru^2\sigma$ sind
$-g_2u^3/12$, $4ru^3$, woraus $4r=g_2/12$ folgt, und demnach erhält
die Differentialgleichung für $\sigma$ die Form:

\weq{10}{\frac{\partial^2 \sigma}{\partial u^2} + 4\eta_1\frac{\partial \sigma}{\partial \omega_1} + 4\eta_2\frac{\partial \sigma}{\partial \omega_2} + \frac{g_2}{12}u^2\sigma = 0,}

und nebenbei ergeben sich aus (9) die Relationen:

\weq{11}{
\begin{aligned}
&\eta_1\frac{\partial \eta_1}{\partial \omega_1} + \eta_2\frac{\partial \eta_1}{\partial \omega_1} = -\frac{g_2}{48}\omega_1, \\ % [sic] The source repeats partial eta_1 / partial omega_1 in the second term; the derivation requires partial eta_1 / partial omega_2.
&\eta_1\frac{\partial \eta_1}{\partial \omega_2} + \eta_2\frac{\partial \eta_2}{\partial \omega_2} = -\frac{g_2}{48}\omega_2. % [sic] The source prints partial eta_1 / partial omega_2 in the first term; the derivation requires partial eta_2 / partial omega_1.
\end{aligned}
}

Nun sollen in der Differentialgleichung (10) statt der Variablen
$\omega_1, \omega_2$ die Variablen $g_2, g_3$ eingeführt werden, und wenn wir also
\[
\begin{aligned}
&\frac{\partial \sigma}{\partial \omega_1} = \frac{\partial \sigma}{\partial g_2}\frac{\partial g_2}{\partial \omega_1} + \frac{\partial \sigma}{\partial g_3}\frac{\partial g_3}{\partial \omega_1}, \\
&\frac{\partial \sigma}{\partial \omega_2} = \frac{\partial \sigma}{\partial g_2}\frac{\partial g_2}{\partial \omega_2} + \frac{\partial \sigma}{\partial g_3}\frac{\partial g_3}{\partial \omega_2}
\end{aligned}
\]
setzen, so erhält (10) die Form:

\weq{12}{\frac{\partial^2 \sigma}{\partial u^2} + h_2\frac{\partial \sigma}{\partial g_2} + h_3\frac{\partial \sigma}{\partial g_3} + \frac{g_2}{12}u^2\sigma = 0,}

worin
\[
\begin{aligned}
h_2 &= 4\left(\eta_1\frac{\partial g_2}{\partial \omega_1} + \eta_2\frac{\partial g_2}{\partial \omega_2}\right), \\
h_3 &= 4\left(\eta_1\frac{\partial g_3}{\partial \omega_1} + \eta_2\frac{\partial g_3}{\partial \omega_2}\right).
\end{aligned}
\]

Die Größen $h_2, h_3$, durch $g_2, g_3$ ausgedrückt, ergeben sich nun
aus der Differentialgleichung selbst, wenn man für $\sigma$ die Entwickelung
 (7) einsetzt. Danach ist bis zur 7ten Potenz von $u$:
%CONT>>
